% GENERATED FILE. Edit canonical lessons, then run npm run generate:books.
% canonical-source-hash: fnv1a64:9b0ca491c4485082
% canonical-lessons: IT-W261-with-a-clock

\chapter{Writing --- With A Clock}
\label{ch:it-with-a-clock}
\hlchaptermodality{261}

\begin{chapteropening}
\textbf{By the end of this chapter:} \emph{I can fill the six fields of a practical Italian form inside a three-minute clock, and I know the shape of this book's twenty-minute A1 writing paper and the CILS rule that one weak part cannot be carried.}

Complete IT-W261-with-a-clock: six form fields against a three-minute timer, then at most one minute checking doubles and accents.
\end{chapteropening}

\section[(with a clock)]{The same writing, with a clock}

\label{lesson:IT-W261-with-a-clock}



\begin{quote}
You wrote your own line with as long as you liked.

A task you can only do untimed is not yet a task you can do.
\end{quote}

\subsection*{Writing — timed production}
\textbf{Whose numbers these are.} This book's A1 writing paper is its own: a project-defined step towards CILS A1, not a copy of it. Its numbers are in \texttt{task-shapes/a1.json}. The paper runs \textbf{20 minutes} and holds \textbf{two} parts: a practical form of \textbf{six to eight} short fields, then a message of \textbf{30 to 40 words} to a named reader, such as \textbf{il nonno}, for a stated purpose. The message has a floor and a ceiling: under thirty words it is not done, and past forty you spend time the task did not ask for. Spelling is marked in both parts, and a held double consonant changes the word: \textbf{nono} is not \textbf{nonno}.

Today you rehearse the first part only, against a short clock. Set a \textbf{three-minute timer}, then fill six fields, in Italian:

\begin{enumerate}
\raggedright
  \item your name;
  \item your surname;
  \item your age with \textbf{ho … anni} (in figures if it is past twenty: this book has spelled numbers only that far);
  \item one person in your family;
  \item one fact about you with \textbf{è};
  \item one more fact you can already write.
\end{enumerate}

When the three-minute timer rings, \textbf{stop}, even mid-word.

\begin{tcolorbox}[breakable,colback=teal!4,colframe=teal!35!black,title={Before you move on}]
Spend at most one minute checking, and do not rewrite the timed answer. Make it one pass for spelling, which the paper marks in both parts. Italian spelling asks two things above all. A \textbf{doubled consonant} is held, and it changes the word: \textbf{nono} is \emph{ninth}, \textbf{nonno} is \emph{grandfather}. And a \textbf{final stressed vowel} carries its accent: \textbf{è}.

Now the rule worth carrying out of this lesson. \textbf{CILS A1}, the real examination this book points at, from the Università per Stranieri di Siena, has \textbf{five} parts: listening, reading, a part on the structures of communication, writing and speaking.

Each part is scored out of \textbf{12}, and you need \textbf{7 in every one}. A strong reading part does not carry a weak writing part. So the part that decides your result is \textbf{the one you are worst at}, not the average.

This book does not repeat how many minutes or words CILS sets for its own writing part. Check the current guide from Siena before you book.
\end{tcolorbox}
